\begin{comment}
%------------------------------------------------------------------------------------------------
\pgfplotstableread[col sep=semicolon, read comma as period]{r_bagging.csv}\datatable

\begin{table}[htbp]
\centering
\caption{RMSEP com RLM para as amostras de a�o-ligas.}
\pgfplotstabletypeset[
    row predicate/.code={%
				%\ifnum#1=0\relax  \pgfplotstableuserowfalse  \fi %ocultando a linha 0
        \ifnum#1=1\relax  \pgfplotstableuserowfalse  \fi %ocultando a linha 1
				\ifnum#1=2\relax  \pgfplotstableuserowfalse  \fi %ocultando a linha 2
    },
    columns={RMSEP, Mn, Mo, Cr, Ni, Fe},
    columns/RMSEP/.style={
        string type,
        string replace*={RLM-bagging}{RLM-\textit{bagging}},
        column name=RMSEP
    },
    columns/Mn/.style={
        fixed, precision=4, column name=Mn,
        postproc cell content/.append code={%
            \ifnum\pgfplotstablerow=1
                \pgfkeysalso{@cell content/.add={}{$^{*17}$}}%
            \fi
        }
    },
    columns/Mo/.style={
        fixed, precision=4, column name=Mo,
        postproc cell content/.append code={%
            \ifnum\pgfplotstablerow=1
                \pgfkeysalso{@cell content/.add={}{$^{*17}$}}%
            \fi
        }
    },
    columns/Cr/.style={
        fixed, precision=4, column name=Cr,
        postproc cell content/.append code={%
            \ifnum\pgfplotstablerow=1
                \pgfkeysalso{@cell content/.add={}{$^{*17}$}}%
            \fi
        }
    },
    columns/Ni/.style={
        fixed, precision=4, column name=Ni,
        postproc cell content/.append code={%
            \ifnum\pgfplotstablerow=1
                \pgfkeysalso{@cell content/.add={}{$^{*17}$}}%
            \fi
        }
    },
    columns/Fe/.style={
        fixed, precision=4, column name=Fe,
        postproc cell content/.append code={%
            \ifnum\pgfplotstablerow=1
                \pgfkeysalso{@cell content/.add={}{$^{*17}$}}%
            \fi
        }
    },
    set decimal separator={,},
    every head row/.style={before row=\toprule, after row=\midrule},
    every last row/.style={after row=\bottomrule}
]{\datatable}
%\legend{$^{*}$ bootstraps utilizados. Fonte:~\oautor}
\legend{Fonte:~\oautor}
\label{tab:rmseprlm}
\end{table}
\end{comment}
%------------------------------------------------------------------------------------------------

\begin{comment}
%-----------------------------------------------------------------------------------------------------
\pgfplotstableread[col sep=semicolon, read comma as period]{r_bagging.csv}\datatable

\begin{table}[htbp]
\centering
\caption{RMSEP da RLM-\textit{bagging}, para cada elemento de a�o-ligas.}
\label{tab:bagging}
\pgfplotstabletypeset[
    row predicate/.code={%
				%\ifnum#1=0\relax  \pgfplotstableuserowfalse  \fi %ocultando a linha 0
        %\ifnum#1=1\relax  \pgfplotstableuserowfalse  \fi %ocultando a linha 1
				\ifnum#1=2\relax  \pgfplotstableuserowfalse  \fi %ocultando a linha 2
    },
    columns={RMSEP, Mn, Mo, Cr, Ni, Fe},
    columns/RMSEP/.style={
        string type,
        string replace*={RLM-bagging}{RLM-\textit{bagging}},
        column name=RMSEP
    },
    columns/Mn/.style={
        fixed, precision=4, column name=Mn,
        postproc cell content/.append code={%
            \ifnum\pgfplotstablerow=1
                \pgfkeysalso{@cell content/.add={}{$^{*17}$}}%
            \fi
        }
    },
    columns/Mo/.style={
        fixed, precision=4, column name=Mo,
        postproc cell content/.append code={%
            \ifnum\pgfplotstablerow=1
                \pgfkeysalso{@cell content/.add={}{$^{*17}$}}%
            \fi
        }
    },
    columns/Cr/.style={
        fixed, precision=4, column name=Cr,
        postproc cell content/.append code={%
            \ifnum\pgfplotstablerow=1
                \pgfkeysalso{@cell content/.add={}{$^{*17}$}}%
            \fi
        }
    },
    columns/Ni/.style={
        fixed, precision=4, column name=Ni,
        postproc cell content/.append code={%
            \ifnum\pgfplotstablerow=1
                \pgfkeysalso{@cell content/.add={}{$^{*17}$}}%
            \fi
        }
    },
    columns/Fe/.style={
        fixed, precision=4, column name=Fe,
        postproc cell content/.append code={%
            \ifnum\pgfplotstablerow=1
                \pgfkeysalso{@cell content/.add={}{$^{*17}$}}%
            \fi
        }
    },
    set decimal separator={,},
    every head row/.style={before row=\toprule, after row=\midrule},
    every last row/.style={after row=\bottomrule}
]{\datatable}
\legend{$^{*}$ bootstraps utilizados. Fonte:~\oautor}
\end{table}
\end{comment}
%--------------------------------------------------------------------
\begin{table}[htbp]
\centering
\caption{RMSEP da RLM e RLM-\textit{subagging} para as amostras de a�o-ligas.}
\label{tab:subagging}
\pgfplotstabletypeset[
    row predicate/.code={%
				%\ifnum#1=0\relax  \pgfplotstableuserowfalse  \fi %ocultando a linha 0
        \ifnum#1=1\relax  \pgfplotstableuserowfalse  \fi %ocultando a linha 1
				%\ifnum#1=2\relax  \pgfplotstableuserowfalse  \fi %ocultando a linha 2
    },
    columns={RMSEP, Mn, Mo, Cr, Ni, Fe},
    columns/RMSEP/.style={
        string type,
        string replace*={RLM-bagging}{RLM-\textit{bagging}},
				string replace*={RLM-subagging}{RLM-\textit{subagging}},
        column name=RMSEP
    },
    columns/Mn/.style={
        fixed, precision=4, column name=Mn,
        postproc cell content/.append code={%
            \ifnum\pgfplotstablerow=2  \pgfkeysalso{@cell content/.add={}{$^{*5}$}}  \fi %exibindo o asterisco na linha 2
        }
    },
    columns/Mo/.style={
        fixed, precision=4, column name=Mo,
        postproc cell content/.append code={%
            \ifnum\pgfplotstablerow=2  \pgfkeysalso{@cell content/.add={}{$^{*16}$}}  \fi %exibindo o asterisco na linha 2
        }
    },
    columns/Cr/.style={
        fixed, precision=4, column name=Cr,
        postproc cell content/.append code={%
            \ifnum\pgfplotstablerow=2  \pgfkeysalso{@cell content/.add={}{$^{*12}$}}  \fi %exibindo o asterisco na linha 2
        }
    },
    columns/Ni/.style={
        fixed, precision=4, column name=Ni,
        postproc cell content/.append code={%
            \ifnum\pgfplotstablerow=2  \pgfkeysalso{@cell content/.add={}{$^{*10}$}}  \fi %exibindo o asterisco na linha 2
        }
    },
    columns/Fe/.style={
        fixed, precision=4, column name=Fe,
        postproc cell content/.append code={%
            \ifnum\pgfplotstablerow=2  \pgfkeysalso{@cell content/.add={}{$^{*8}$}}  \fi %exibindo o asterisco na linha 2
        }
    },
    set decimal separator={,},
    every head row/.style={before row=\toprule, after row=\midrule},
    every last row/.style={after row=\bottomrule}
]{\datatable}
\legend{$^{*}$ bootstraps utilizados. Fonte:~\oautor}
\end{table}
%-----------------------------------------------------------------------------------

\begin{comment}
	
\pgfplotstableread[col sep=comma]{r_bagging.csv}\datatable

\begin{table}[htbp]
\centering
\caption{Compara��o de Desempenho}
\pgfplotstabletypeset[
    columns={RMSEP, Mn, Mo, Cr, Ni, Fe},
    columns/RMSEP/.style={string type, column name=\textbf{RMSEP}},
    columns/Mn/.style={fixed, precision=4, column name=\textbf{Mn}},
    columns/Mo/.style={fixed, precision=4, column name=\textbf{Mo}},
    columns/Cr/.style={fixed, precision=4, column name=\textbf{Cr}},
    columns/Ni/.style={fixed, precision=4, column name=\textbf{Ni}},
		columns/Fe/.style={fixed, precision=4, column name=\textbf{Fe}},
    every head row/.style={before row=\toprule, after row=\midrule},
    every last row/.style={after row=\bottomrule},
]{\datatable}
\legend{Fonte:~\oautor}
\end{table}
\end{comment}

%-------------------------------------------------------------------------------------

\pgfplotstableread[col sep=semicolon, read comma as period]{r_bagging.csv}\datatable

\begin{table}[htbp]
\centering
\caption{RMSEP da RLM-bagging, para cada elemento de a�o-ligas.}
\pgfplotstabletypeset[
    columns={RMSEP, Mn, Mo, Cr, Ni, Fe},
    columns/RMSEP/.style={
        string type,
        string replace*={bagging}{\textit{bagging}},
        column name=\textbf{RMSEP}
    },
    columns/Mn/.style={fixed, precision=4, column name=\textbf{Mn},
                       dec sep align},
    columns/Mo/.style={fixed, precision=4, column name=\textbf{Mo},
                       dec sep align},
    columns/Cr/.style={fixed, precision=4, column name=\textbf{Cr},
                       dec sep align},
    columns/Ni/.style={fixed, precision=4, column name=\textbf{Ni},
                       dec sep align},
    columns/Fe/.style={fixed, precision=4, column name=\textbf{Fe},
                       dec sep align},
    set decimal separator={,},
    every head row/.style={before row=\toprule, after row=\midrule},
    every last row/.style={after row=\bottomrule}
]{\datatable}
\legend{Fonte:~\oautor}
\end{table}

%-------------------------------------------------------------------------------------


%-----------------------------------------------------------------------
\pgfplotstableread[col sep=semicolon, read comma as period]{r_bagging.csv}\datatable

\begin{table}[htbp]
\centering
\caption{RMSEP da RLM-\textit{bagging}, para cada elemento de a�o-ligas.}
\pgfplotstabletypeset[
    columns={RMSEP, Mn, Mo, Cr, Ni, Fe},
    columns/RMSEP/.style={
        string type,
        string replace*={RLM-bagging}{RLM-\textit{bagging}},
        column name=RMSEP
    },
    columns/Mn/.style={
        fixed, precision=4, column name=Mn,
        postproc cell content/.append code={%
            \ifnum\pgfplotstablerow=1
                \pgfkeysalso{@cell content/.add={}{$^{*17}$}}%
            \fi
        }
    },
    columns/Mo/.style={
        fixed, precision=4, column name=Mo,
        postproc cell content/.append code={%
            \ifnum\pgfplotstablerow=1
                \pgfkeysalso{@cell content/.add={}{$^{*17}$}}%
            \fi
        }
    },
    columns/Cr/.style={
        fixed, precision=4, column name=Cr,
        postproc cell content/.append code={%
            \ifnum\pgfplotstablerow=1
                \pgfkeysalso{@cell content/.add={}{$^{*17}$}}%
            \fi
        }
    },
    columns/Ni/.style={
        fixed, precision=4, column name=Ni,
        postproc cell content/.append code={%
            \ifnum\pgfplotstablerow=1
                \pgfkeysalso{@cell content/.add={}{$^{*17}$}}%
            \fi
        }
    },
    columns/Fe/.style={
        fixed, precision=4, column name=Fe,
        postproc cell content/.append code={%
            \ifnum\pgfplotstablerow=1
                \pgfkeysalso{@cell content/.add={}{$^{*17}$}}%
            \fi
        }
    },
    set decimal separator={,},
    every head row/.style={before row=\toprule, after row=\midrule},
    every last row/.style={after row=\bottomrule}
]{\datatable}
\legend{$^{*}$ 17 bootstraps utilizados. Fonte:~\oautor}
\end{table}

%---------------------------------------------------------
\pgfplotstableread[col sep=semicolon, read comma as period]{r_bagging.csv}\datatable

\begin{table}[htbp]
\centering
\caption{RMSEP da RLM-bagging, para cada elemento de a�o-ligas.}
\pgfplotstabletypeset[
    row predicate/.code={%
				%\ifnum#1=0\relax  \pgfplotstableuserowfalse  \fi %ocultando a linha 0
        %\ifnum#1=1\relax  \pgfplotstableuserowfalse  \fi %ocultando a linha 1
				%\ifnum#1=2\relax  \pgfplotstableuserowfalse  \fi %ocultando a linha 2
    },
    columns={RMSEP, Mn, Mo, Cr, Ni, Fe},
    columns/RMSEP/.style={
        string type,
        string replace*={RLM-bagging}{RLM-\textit{bagging}},
				string replace*={RLM-subagging}{RLM-\textit{subagging}},
        column name=RMSEP
    },
    columns/Mn/.style={
        fixed, precision=4, column name=Mn,
        postproc cell content/.append code={%
            %\ifnum\pgfplotstablerow=1  \pgfkeysalso{@cell content/.add={}{$^{*}$}}  \fi %exibindo o asterisco na linha 1
        }
    },
    columns/Mo/.style={
        fixed, precision=4, column name=Mo,
        postproc cell content/.append code={%
            %\ifnum\pgfplotstablerow=1  \pgfkeysalso{@cell content/.add={}{$^{*}$}}  \fi %exibindo o asterisco na linha 1
        }
    },
    columns/Cr/.style={
        fixed, precision=4, column name=Cr,
        postproc cell content/.append code={%
            %\ifnum\pgfplotstablerow=1  \pgfkeysalso{@cell content/.add={}{$^{*}$}}  \fi %exibindo o asterisco na linha 1
        }
    },
    columns/Ni/.style={
        fixed, precision=4, column name=Ni,
        postproc cell content/.append code={%
            %\ifnum\pgfplotstablerow=1  \pgfkeysalso{@cell content/.add={}{$^{*}$}}  \fi %exibindo o asterisco na linha 1
        }
    },
    columns/Fe/.style={
        fixed, precision=4, column name=Fe,
        postproc cell content/.append code={%
            %\ifnum\pgfplotstablerow=1  \pgfkeysalso{@cell content/.add={}{$^{*}$}}  \fi %exibindo o asterisco na linha 1
        }
    },
    set decimal separator={,},
    every head row/.style={before row=\toprule, after row=\midrule},
    every last row/.style={after row=\bottomrule}
]{\datatable}
\legend{$^{*}$ bootstraps utilizados. Fonte:~\oautor}
\end{table}

%------------------------------------------------------------

\begin{table}[htbp]
\centering
\caption{Desempenho comparativo das t�cnicas de calibra��o multivariada para elementos qu�micos em a�o-liga.}
\label{tab:baggingsubba}
\footnotesize
\begin{tabular}{clccccl}
\toprule
\textbf{Elemento} & \textbf{T�cnica} & \textbf{RMSEP} & \textbf{MAE} & 
\textbf{R�} & \textbf{Tempo} & \textbf{Par�metros} \\
\midrule
\multirow{3}{*}{Mn} 
  & RLM              & 0,1051 & 0,0887 & 0,9907 &  2,4 ms & � \\
  & RLM-\textit{bagging}      & 0,1072 & 0,0908 & 0,9903 & 68,2 ms & $m^* = 30$ \\
  & RLM-\textit{subagging}$^*$ & \textbf{0,0966} & \textbf{0,0805} & \textbf{0,9922} & 17,1 ms & $k^*=80\%, m^*=10$ \\
\midrule
\multirow{3}{*}{Mo} 
  & RLM$^*$          & \textbf{0,2633} & \textbf{0,1952} & \textbf{0,9882} &  2,4 ms & � \\
  & RLM-\textit{bagging}      & 0,3388 & 0,2722 & 0,9805 & 67,5 ms & $m^* = 30$ \\
  & RLM-\textit{subagging}    & 0,2879 & 0,2201 & 0,9859 & 46,8 ms & $k^*=80\%, m^*=30$ \\
\midrule
\multirow{3}{*}{Cr} 
  & RLM$^*$          & \textbf{1,0798} & \textbf{0,8654} & \textbf{0,9938} &  2,5 ms & � \\
  & RLM-\textit{bagging}      & 1,2089 & 0,9830 & 0,9922 & 44,5 ms & $m^* = 20$ \\
  & RLM-\textit{subagging}    & 1,1294 & 0,9010 & 0,9932 & 46,8 ms & $k^*=80\%, m^*=30$ \\
\midrule
\multirow{3}{*}{Ni} 
  & RLM              & 1,6043 & 1,3256 & 0,9837 &  2,3 ms & � \\
  & RLM-\textit{subagging}$^*$ & \textbf{1,6003} & \textbf{1,2409} & \textbf{0,9838} & 10,6 ms & $k^*=60\%, m^*=10$ \\
  & RLM-\textit{bagging}      & 1,7085 & 1,4122 & 0,9815 & 66,6 ms & $m^* = 30$ \\
\midrule
\multirow{3}{*}{Fe} 
  & RLM              & 2,7513 & 2,1865 & 0,9892 & 13,5 ms & � \\
  & RLM-\textit{bagging}$^*$   & \textbf{2,6406} & \textbf{2,0246} & \textbf{0,9900} & 97,2 ms & $m^* = 30$ \\
  & RLM-\textit{subagging}    & 2,6509 & 2,0578 & 0,9900 & 41,2 ms & $k^*=80\%, m^*=30$ \\
\bottomrule
\end{tabular}
\begin{flushleft}
\footnotesize
$^*$ Melhor desempenho em RMSEP. Par�metros $m^*$ (n�mero de replica��es 
\textit{bootstrap}) e $k^*$ (fra��o de subamostra) otimizados por 
valida��o cruzada 5-\textit{fold} aplicada exclusivamente ao conjunto 
de calibra��o (\textit{seed} = 42). O conjunto de teste foi utilizado 
uma �nica vez para avalia��o final.
\end{flushleft}
\legend{Fonte:~\oautor}
\end{table}
%-----------------------------------------------------------

\begin{table}[!htbp]
\centering
\caption{O coeficiente de varia��o.}
\begin{tabular}{lllllll}
\hline
 &  & $m \to$ &  &  &  &  \\ 
 &  & 25 & 50 & 100 & 200 & $\infty$ \\ 
\hline
$cv(\widehat{se}_{\infty})$ & 0.25 & 0.29 & 0.27 & 0.26 & 0.25 & 0.25 \\ 
\multicolumn{1}{c}{$\downarrow$} & 0.20 & 0.24 & 0.22 & 0.21 & 0.21 & 0.20 \\ 
 & 0.15 & 0.21 & 0.18 & 0.17 & 0.16 & 0.15 \\ 
 & 0.10 & 0.17 & 0.14 & 0.12 & 0.11 & 0.10 \\ 
 & 0.05 & 0.15 & 0.11 & 0.09 & 0.07 & 0.05 \\ 
 & 0.00 & 0.14 & 0.10 & 0.07 & 0.05 & 0.00 \\ 
\hline
\end{tabular}\\[0.2ex]
\legend{Fonte:~An Introduction to the Bootstrap. USA: Chapman \& Hall/CRC, 1993.}
\label{tab:coefvaria}
\end{table}